\subsection{Empirical Distribution and Survival Summaries}

Exploratory summaries first treat the observed values as a batch of numbers.  A stochastic model may later explain their variability, but the displays themselves are useful before a model has been chosen.

\begin{definition} \label{def:empirical-cdf-and-survival-summary} \hlt{Empirical CDF and Empirical Survival Function}\\
Let $x_1,\ldots,x_n\in\R$ be data; for survival summaries assume $x_i\geq0$ for $1\leq i\leq n$.
\begin{enumerate}[label=(\roman*)]
\setlength{\itemsep}{0pt}
\item \hlt{Empirical CDF}: $F_n:\R\to[0,1]$, $F_n(x)\defeq n^{-1}\#\{i\in\{1,\ldots,n\}:x_i\leq x\}$.
\item \hlt{Empirical Survival Function}: $S_n:[0,\infty)\to[0,1]$, $S_n(t)\defeq1-F_n(t)=n^{-1}\#\{i\in\{1,\ldots,n\}:x_i>t\}$.
\end{enumerate}
\hlt{Key Insight}: $F_n$ and $S_n$ encode complementary empirical tails, with $S_n(t)=1-F_n(t)$ for $t\geq0$.
\end{definition}

\begin{remark} \label{rem:ecdf-step-convention} \hlt{Step-Function Convention}\\
Let $v_1<\cdots<v_m$ be the distinct observed values, with multiplicities $r_1,\ldots,r_m$. Then $F_n$ is right-continuous, constant on each component of $\R\setminus\{v_1,\ldots,v_m\}$, and has jump $\Delta F_n(v_j)=\frac{r_j}{n}$.
\end{remark}

\begin{figure}[H]
  \centering
  % Source: supplied Chapter 10 screenshots on empirical CDF and survival displays.
% Schematic redraw, not a reproduction of the scanned figures.
\begin{tikzpicture}[
  x=0.9cm,
  y=2.8cm,
  >=Stealth,
  line cap=round,
  line join=round,
  every node/.style={font=\scriptsize},
  axis/.style={->, gray!70, line width=0.45pt},
  tick/.style={gray!70, line width=0.35pt},
  ecdf/.style={blue!70!black, line width=0.95pt},
  surv/.style={green!45!black, line width=0.95pt},
  guide/.style={gray!35, dashed, line width=0.35pt}
]
  \begin{scope}
    \draw[axis] (0,0) -- (4.75,0) node[black!78,right] {$x$};
    \draw[axis] (0,0) -- (0,1.12) node[black!78,above] {$F_n(x)$};
    \foreach \x in {.55,1.25,2.05,3.0,3.75}
      \draw[tick] (\x,0.015) -- (\x,-0.015);
    \foreach \y/\lab in {0/0,.5/.5,1/1}
      \draw[tick] (0.04,\y) -- (-0.04,\y) node[black!78,left] {$\lab$};
    \draw[guide] (0,.5) -- (4.45,.5);
    \draw[ecdf]
      (0,0) -- (.55,0) -- (.55,.2) -- (1.25,.2) -- (1.25,.4) --
      (2.05,.4) -- (2.05,.6) -- (3.0,.6) -- (3.0,.8) --
      (3.75,.8) -- (3.75,1.0) -- (4.45,1.0);
  \end{scope}

  \begin{scope}[xshift=5.8cm]
    \draw[axis] (0,0) -- (4.75,0) node[black!78,right] {$t$};
    \draw[axis] (0,0) -- (0,1.12) node[black!78,above] {$S_n(t)$};
    \foreach \x in {.55,1.25,2.05,3.0,3.75}
      \draw[tick] (\x,0.015) -- (\x,-0.015);
    \foreach \y/\lab in {0/0,.5/.5,1/1}
      \draw[tick] (0.04,\y) -- (-0.04,\y) node[black!78,left] {$\lab$};
    \draw[guide] (0,.5) -- (4.45,.5);
    \draw[surv]
      (0,1.0) -- (.55,1.0) -- (.55,.8) -- (1.25,.8) -- (1.25,.6) --
      (2.05,.6) -- (2.05,.4) -- (3.0,.4) -- (3.0,.2) --
      (3.75,.2) -- (3.75,0) -- (4.45,0);
  \end{scope}
\end{tikzpicture}

  \caption{Empirical CDF (blue) rises as empirical survival (green) falls.}
  \label{fig:eda-ecdf-survival-steps}
\end{figure}

\begin{method} \label{meth:empirical-log-survival-plot} \hlt{Empirical Log-Survival Plot}\\
Suppose $0\leq T_{(1)}<\cdots<T_{(n)}$ and there are no ties. Set $p_i\defeq\frac{i}{n+1}$ and $y_i\defeq\log(1-p_i)$.
\begin{enumerate}[label=(\roman*)]
\setlength{\itemsep}{0pt}
\item \hlt{Plot}: Draw $P_{\log S}^{(n)}\defeq\{(T_{(i)},y_i):1\leq i\leq n\}$; using $n+1$ keeps $y_n$ finite.
\item \hlt{Hazard Reading}: Since $(\log S)'=-h$, more negative local slopes indicate larger hazards.
\item \hlt{Tail Variability}: $\operatorname{Var}(\log[1-F_n(t)])\approx n^{-1}\frac{F(t)}{1-F(t)}$, which grows as $F(t)\uparrow1$.
\end{enumerate}
\end{method}
